%% kapa: exam=1 round=27 year=2016 period=1 form=A subject=real_estate subject_ko=부동산학원론 seq=34 number=114
%% source: Q-Net 감정평가사 기출문제 (cst00309) artlSeq=5169402 file=extracted/1008/5169402/제27회 감정평가사 1차 문제/제27회 감정평가사 1교시문제(A형).pdf
%% conversion: kordoc 4.24.51 -> tools/gen1.py (auto); review=proofread-sonnet; flags=none
\begin{question}{114}
부동산정책에 관한 설명으로 옳지 않은 것은?
\begin{choices}
  \item 부동산정책이란 바람직한 부동산활동을 유도하기 위한 목표설정과 이를 달성하기 위한 각종 부동산대책의 결정 및 운용에 관한 정부의 공적인 계획이나 실행 행위를 말한다.
  \item 부동산 거래신고 제도는 부동산 거래신고에 관한 법령에 따라 거래당사자가 부동산등에 관한 매매계약을 체결한 경우 그 실제 매매가격 등을 거래계약 후 잔금일로부터 60일 이내에 그 부동산등의 소재지를 관할하는 시장ㆍ군수 또는 구청장에게 공동 또는 예외적인 경우 단독으로 신고하게 하여 건전하고 투명한 부동산 거래질서를 확립하여 국민경제에 이바지함을 목적으로 한다.
  \item 개발제한구역의 지정 및 관리에 관한 특별조치법령상 국토교통부장관은 국방부장관의 요청으로 보안상 도시의 개발을 제한할 필요가 있다고 인정되면 개발제한구역의 지정 및 해제를 도시ㆍ군관리계획으로 결정할 수 있다.
  \item 지적재조사사업은 공간정보의 구축 및 관리 등에 관한 법령에 따라 지적공부의 등록사항을 조사ㆍ측량하여 기존의 지적공부를 디지털에 의한 새로운 지적공부로 대체함과 동시에 지적공부의 등록사항이 토지의 실제 현황과 일치하지 아니하는 경우 이를 바로 잡기 위하여 실시하는 국가사업으로 국토를 효율적으로 관리함과 아울러 국민의 재산권 보호에 기여함을 목적으로 한다.
  \item 산지관리법령상 국가나 지방자치단체는 산지전용ㆍ일시사용제한지역의 지정목적을 달성하기 위하여 필요하면 산지소유자와 협의하여 산지전용ㆍ일시사용제한지역의 산지를 매수할 수 있다.
\end{choices}
\end{question}
\answer{2}
